\section{Codeforces 刷题总结}
%\textcolor{blue}{\textbf{271A}}
%\begin{lstlisting}[language=C++]\end{lstlisting}
在本节中，我将总结我在 Codeforces 上做题的经验。

\subsection{遇到的问题以及学到的方法}
\begin{itemize}
	\subsection{完美的一年}
	\item 题目链接：\href{https://codeforces.com/problemset/problem/271/A}{217A(点击此处)};方法:在判断一个字符串是否每个字符均不相等时可以用\textcolor{red}{set集合来去重}，\textcolor{red}{最终比较set的长度和原字符串的长度}来判断是非满足条件；
	
	代码示例：\\
\begin{lstlisting}[language=C++]
void solve() {
	ll n, num;
	cin >> n;
	num = n + 1;
	while (num) {
		string s;
		s = to_string(num);
		set<char> st(s.begin(), s.end());
		if (s.size() == st.size()) {
			cout << s << endl;
			break;
		}
		num++;
	}
}
\end{lstlisting}
\subsection{学校排队(与LuoguP3741类似)}
\item 题目链接:\href{https://codeforces.com/problemset/problem/266/B}{266B学校排队(点击此处)}；总结：通过字符串交换来达到目的的题都可以\textcolor{red}{用X填充思想；}
	
代码示例：\\
\begin{lstlisting}[language=C++]
void solve() {
	ll n, t;
	cin >> n >> t;
	string s1, s2;
	cin >> s1;
	s2 = s1;
	while (t--) {
		for (ll i = 0; i < n - 1; i++) {
			if (s1[i] == 'B' && s1[i + 1] == 'G' && s2[i] != 'X') {
				swap(s1[i], s1[i+1]);
				s2[i] = 'X';
				s2[i + 1] = 'X';
			}
		}
		s2=s1;
	}
	cout << s1 << endl;
}
\end{lstlisting}
\subsection{向上取整问题，不用ceil函数}
\item 题目链接：\href{https://codeforces.com/contest/1328/problem/A}{629A整除性问题(点击此处)}；

代码示例：\\
\begin{lstlisting}[language=C++]
void solve() {
	ll a, b;
	cin >> a >> b;
	if (a % b == 0) {
		cout << 0 << endl;
	} else {
		ll c = (a + b - 1) / b;//代替ceil函数实现向上取整即(a+b-1)/b
		ll d = c * b;
		cout << d - a << endl;
	}
}
\end{lstlisting}
\end{itemize}
\newpage
% \subsection{代码示例}
% 下面是一个代码示例：

% \begin{lstlisting}[language=C++]
	% #include <bits/stdc++.h>
	% using namespace std;
	
	% int main() {
		%     cout << "Hello Codeforces!" << endl;
		%     return 0;
		% }
	% \end{lstlisting}
